\section{Types and static semantics}
\label{sec:types}

Checking is static and \emph{fatal}: a type error is a compile error.
The type system is small --- scalars, arrays, records, maps, options,
pointers, functions and \ic{void} --- and mostly structural.

\subsection{Scalar types}

\begin{center}
\begin{tabular}{lll}
\toprule
type & domain & notes \\
\midrule
\ic{long}  & 64-bit signed integers & wraparound on overflow \\
\ic{byte}  & $0\ldots255$ & stored in one byte inside \ic{[byte]} \\
\ic{bool}  & \ic{true}, \ic{false} & \\
\ic{float} & IEEE-754 binary64 & \\
\ic{str}   & immutable byte strings & length-prefixed; may hold NUL \\
\bottomrule
\end{tabular}
\end{center}

\subsection{Arrays}
\label{sec:type-arrays}

\ic{[T]} is a heap array of \ty{T}, \ic{[T; N]} a fixed-length array of
\ty{T} with \ty{N} a literal. The element type is part of the type and
arrays are \textbf{homogeneous}: every element of an array has exactly
the element type. Heterogeneous array literals are rejected:

\begin{lstlisting}
let bad = [1, "two"];        // error: array element: expected long, got str
\end{lstlisting}

\noindent This applies recursively --- \ic{[[1], [2]]} has type
\ic{[[long]]}, and every operator that mixes arrays and elements
(Section~\ref{sec:arrays}) enforces the same rule.

\subsection{Records}
\label{sec:type-records}

A record literal \ic{\{ a: 1, b: "x" \}} has the structural shape
\ic{\{ a: long, b: str \}}; field order is preserved and is part of the
type. A field read is legal when the static shape has the field. The
shape-hole form \ic{\{ k: str \}} matches any record with at least those
fields and is used where the precise shape is not known statically (it
poisons quietly rather than erroring --- see
Section~\ref{sec:poison}).

\subsection{Maps}

\ic{\{K => V\}} is a dynamic map from hashable keys \ty{K} to one
value type \ty{V}. Reading \ic{m[k]} always yields \ic{V?}; see
Section~\ref{sec:maps}.

\ty{K} is a hashable key type (\ic{str}, \ic{long}, \ic{byte},
\ic{bool}, \ic{float}), but a \textbf{map literal only accepts \ic{str}
keys} --- the checker's inference for a literal fixes the key type to
\ic{str}, so \ic{\{ 1: 5 \}} infers a shape rather than a map. Build a
non-\ic{str}-keyed map from an annotated empty literal and store into
it (Section~\ref{sec:maps}).

\subsection{Options and pointers}
\label{sec:type-options}

\ic{T?} is ``a \ty{T} or nothing''. Its empty case is \ic{none};
\ic{some}/\ic{unwrap}/\ic{is\_some} and the postfix unwrap \ic{x!}
(Section~\ref{sec:builtin-option}) manipulate it. There is no
\ic{null}: \ic{none} is the only empty value, and it exists only at
option types --- a plain \ic{str} or \ic{long} can never be empty, so
\ic{s == none} on a \ic{str} is a type error. On its own \ic{none} has
the type ``option of anything'' (printed \ic{none}) and adopts the
option type its context expects. A plain \ty{T} does not widen to
\ic{T?}; write \ic{some(v)}. \ic{*T} is a raw machine pointer; \ic{\&e}
yields one (Section~\ref{sec:pointers}).

\subsection{Void}

\ic{void} is the type of an expression that computes no value: a loop,
an \ic{if} without \ic{else}, a call to a void function
(Section~\ref{sec:void}), a bare \ic{return}. Void is not a value: it
cannot be bound, passed, printed or returned as one.

\paragraph{Disagreeing \ic{if} arms.} When the two arms of an
\ic{if}/\ic{else} have different types, the whole \ic{if} is
\ic{void} --- a statement. That is legal wherever a statement is, and
an error wherever a value is needed:

\begin{lstlisting}
if (ok) { 1 } else { "one" };            // fine: a statement
let v = if (ok) { 1 } else { 2 };        // long
let w = if (ok) { 1 } else { "one" };    // error: cannot bind a void value to 'w'
\end{lstlisting}

\subsection{Functions}

\ic{(T, T, ...) -> R}. Return types are inferred when omitted and
checked when written. A parameter annotation is optional: without one
the parameter takes the unknown type, which poisons quietly
(Section~\ref{sec:poison}) and is treated as a machine word by the
backend --- so annotate every parameter that carries a string, array,
record or option. Function values are first-class (they may be bound,
passed, stored and called); see Section~\ref{sec:functions}.

\subsection{Annotations and inference}

\begin{itemize}[nosep]
\item \ic{def} parameter annotations are \textbf{recommended}; a missing
  one takes the unknown type and never cascades a second diagnostic
  (Section~\ref{sec:poison}).
\item Return types are inferred from the body when omitted; when written,
  every \ic{return} and the body's trailing value must conform.
\item \ic{let} annotations are optional and checked; an annotated
  \emph{array literal} \emph{adopts} the annotated element type:
  \ic{let b: [byte] = [104, 105];} packs those elements one byte each.
  Adoption is restricted to literals --- assigning an existing
  \ic{[long]} to a \ic{[byte]} binding is a type error (the layouts
  differ).
\item An empty array literal \ic{[]} is legal unannotated: there is
  nothing to infer an element type from, so the element type is
  \emph{unknown} and it poisons quietly (Section~\ref{sec:poison}).
  Annotate it to pin the element type --- \ic{let b: [byte] = [];}.
  An unannotated \ic{\{\}} is the empty record; under a map annotation
  it is the empty map: \ic{let m: \{str => long\} = \{\};}.
\end{itemize}

\subsection{Conversions}
\label{sec:conversions}

The implicit conversions are:
\ic{long $\bs to$ float} (the only numeric widening),
\ic{byte $\bs to$ long}, \ic{byte $\bs to$ float}, and the narrowing
\ic{long $\bs to$ byte} at a \ic{byte} position (the value is masked to
eight bits: \ic{let b: byte = 258;} stores 2). A \ic{str} also
flows into a \ic{*T} position, since a string value is a byte address.
Apart from these, nothing converts: there is no cast syntax.

A value converts where it is stored under a known type: an annotated
\ic{let}; an assignment to a variable, element, field, map entry or
pointee; an argument; a \ic{return} or trailing value under a declared
return type; an element of an array literal; and the element that
\ic{+=}, \ic{*=} or \ic{::} adds. Array \ic{+} and \ic{+=} with an
array operand convert element by element into the left operand's
element type: \ic{[0.5] + [7, 8]} is \ic{[0.5, 7.0, 8.0]}. Arithmetic
on a \ic{byte} is \ic{long} arithmetic (\ic{-b} is a negative
\ic{long}); storing the result back into a \ic{byte} masks it.

Nothing converts a value that is already inside another one. Arrays,
maps and records are references and an option is a box, so an existing
\ic{[long]} is not a \ic{[float]}, and a record \ic{\{x: long\}} is not
a \ic{\{x: float\}}: a record's numeric field types must match exactly.

\subsection{Operator typing}

Arithmetic requires numbers; if either side is \ic{float} the result is
\ic{float} and the other side widens. \ic{+} additionally allows
strings (concatenation, with auto-stringification of the non-string
side) and arrays (Section~\ref{sec:arrays}); arrays win over strings
when both appear. Comparisons yield \ic{bool}; \ic{<}\,\ldots\ic{>=}
on two strings is lexicographic byte order. \ic{\&\&}/\ic{||} accept any
values: a non-zero value is true and the result is \ic{bool}
(Section~\ref{sec:evaluation}). Equality works on any two values of the
same type; it is by content for scalars and strings, structural for
arrays, and by identity for records and maps (Section~\ref{sec:equality}).

\subsection{Compound assignment typing}
\label{sec:opassign-typing}

For a target of type \ty{T} and value of type \ty{U}:

\begin{center}
\begin{tabular}{lll}
\toprule
\ty{T} & operators & rule \\
\midrule
numbers & \ic{+= -= *= /= \%=} & \ty{U} converts to \ty{T}; stored as \ty{T} \\
\ic{str} & \ic{+=} only & \ty{U} $=$ \ty{str} \\
array & \ic{+=}, \ic{*=} only & Section~\ref{sec:arrays} \\
\bottomrule
\end{tabular}
\end{center}

\noindent Anything else is a compile error, e.g. ``\ic{arrays support
only '+=' and '*=' (in place), not '-='}''. \ic{++}/\ic{--} require a
numeric target.

\subsection{The poison law}
\label{sec:poison}

Missing information (a missing annotation, an unanalyzable expression)
produces an \emph{unknown} type that unifies with everything. The law
this enables: \textbf{the root cause is reported exactly once} --- an
unknown-typed piece never generates follow-on errors at every enclosing
use. Concretely, unknown unifies silently in element checks, call
arguments and branches (so an unknown arm never makes an \ic{if}
void); the missing annotation itself is the single diagnostic.
